\section{Experimental Design for the Extensions}\label{sec:methodology}
\subsection{What Is Tested, and on Which Programs?}
The evaluation separates three practical questions: whether Zone identifies
embedded data correctly, whether a transformed program still performs its
expected task, and how often static scanners flag the resulting file.
Table~\ref{tab:test_populations} summarizes the inputs used for each question.

\begin{table}[t]
\caption{Test populations. The 13 controlled programs and the 13
selected x64 payloads are different sets with different purposes.}
\label{tab:test_populations}
\centering\small
\begin{tabularx}{\linewidth}{@{}l r Y@{}}
\toprule
Population & Inputs & Purpose \\
\midrule
Controlled functional programs & 13 & Check exact expected output after transformation. \\
Annotated classification programs & 6 & Check which bytes are identified as data. \\
Original Windows payload corpus & 102 & Measure construction, functionality, and static detection. \\
Recovery subset of that corpus & 49 & Compare historical and updated working rates on previously problematic inputs. \\
Data-handling subset of that corpus & 19 & Compare symbolic and LLM identification within the same pipeline. \\
\bottomrule
\end{tabularx}
\end{table}

The 13 \emph{controlled functional programs} (test fixtures) are purpose-built
assembly programs of 19--68 bytes. Six per architecture cover arithmetic,
branches, loops, local calls, combined control flow, and embedded strings;
one extra x64 case accesses data through a computed code address. Each
writes a known result into a 256-byte buffer. A native C runner compares
the entire buffer, including untouched positions, with the expected bytes.
All 13 originals pass.

The six \emph{annotated classification programs} instead test data
locations: two contain multiple strings, one a binary table, two no
separate data, and one instruction bytes resembling text. Correct regions
are annotated before either classifier runs and withheld from the model.
These suites check individual mechanisms; original payloads provide the
broader application-level evaluation.

The original corpus contains 38 x64 and 64 x86 Windows payloads.
Functional tests check the expected dialog, network connection event,
or clean termination for the relevant payload class. All 102 originals pass.
Throughout this article, \emph{working} means that the transformed output
passes the corresponding functional test. We use this empirical definition
consistently for the historical and extended evaluations.

The recovery comparison selects the 49 original inputs whose historical
working rate was below 50\%: 13 x64 and 36 x86. It compares the same
inputs and 15 transformation configurations before and after the
pipeline update. New results use ten variants per x64 input/configuration
and five per x86 input/configuration. Data-preparation-only controls are
reported separately. The comparison measures the updated pipeline as a
whole, including the improved MessageBox test; it is not an ablation of
data identification alone.

\subsection{Configurations, Repetitions, and Model Calls}
A \emph{cell} combines an input, transformation, and prompting condition:
for example, one x64 program transformed with EQ and two examples.
Repeating that combination gives repetitions of the same cell.

Both LLM roles request temperature 0.2 without a sampling seed or top-p;
these client settings do not imply identical repeated answers.

The native controlled benchmark uses three runs per input and transformation
family. The LLM benchmark uses four families, two prompting conditions,
and three separate single-generation requests per cell. This gives
$7\times4\times2\times3=168$ x64 requests and
$6\times4\times2\times3=144$ x86 requests. All rejected responses and
transport failures remain in the counts. The classification benchmark uses
three LLM requests and one symbolic analysis per annotated program.

On the full original corpus, wrapping uses three randomization seeds per
input, giving $102\times3=306$ planned variants. The LLM experiment has
$102\times4\times2=816$ cells and an output-token limit of 65,536.
It targets one accepted output per cell; two successful cells required a
second generation, and fourteen failed cells were restarted without success.
These results describe the completed campaign, including those retries.

The data-handling comparison uses 19 original inputs (13 x64,
six x86) with data extracted by symbolic preparation. Each method proposes
data regions, after which the same deterministic handling and functional
test are applied. Input selection and shared processing are thus held
explicit in interpreting the comparison.

The accompanying artifact records engine versions and request histories.
Native x86 IBP was tested with an earlier implementation than the later
x86 LLM and wrapping experiments; these are separately identified in the
artifact. The native x86 IBP result therefore characterizes that recorded
version; no current-version retest is available, and no version-improvement
claim is inferred from comparing these campaigns.

\subsection{Build Rate, Working Rate, and Classification Accuracy}
Let $N$ be planned cases, $B$ produced binaries, and $W$ binaries passing
the functional test. Build rate is $B/N$, overall working rate is $W/N$,
and working rate among produced binaries is $W/B$.
For example, ten planned cases, eight produced binaries, and six successful
tests give 80\%, 60\%, and 75\%, respectively. The selected-cohort comparison divides
successful tests by recorded tests, matching the available historical
records. The overall update in Table~\ref{tab:overall_working_update}
combines the new rates for retested inputs with the retained historical
rates for other inputs. Each input keeps its historical test-count weight.
This avoids changing the composition of the x86 average merely because
the new campaign uses fewer repetitions. No missing historical outcome is
imputed. The final table row is the arithmetic mean of the 15
configuration rates, not weighted by their sample counts. Detailed
contributions are supplied in Supplementary Material S1. Repeated
executions of one binary are not new variants.

For classification, predicted positions $P$ are bytes labelled as data;
reference positions $G$ are the preannotated correct data bytes. Precision
is $|P\cap G|/|P|$ and recall is $|P\cap G|/|G|$. For a reference string
at bytes 20--29, predicting 20--27 gives 100\% precision and 80\% recall.
Exact matching requires all intervals to agree; programs without separate
data require an empty set. For unannotated original payloads, the evaluation
reports agreement between methods and downstream functionality.

\subsection{Similarity, Detection Rate, Diversity, and Cost}
\paragraph{Representation measurements.}
Historical cosine and entropy definitions are retained. The LLM client's
SequenceMatcher compares prepared and returned listings, filtering
NOP-like operations in the NOP condition. This diagnostic is supplied in
the supplementary records rather than substituted for the historical metrics.

\paragraph{Detection rate.}
We use the same engine-level definition as the historical evaluation:
$d_i=m_i/a_i$, where $m_i$ is the number of malicious responses for file
$i$ and $a_i$ is its total reported engine-response count. Tables report
$\overline d$, the mean across files, as a percentage. For example,
3 malicious responses among 75 give a detection rate of 4\%.
Nonempty current reports contain 75 responses. Unsupported formats,
timeouts, and failed analyses remain in that denominator, consistently
with the historical convention; their breakdown and report-completeness
checks are included in the supplementary records. Missing reports are not
assigned a zero detection rate. Each distinct output hash is scanned once.

\paragraph{Diversity and resources.}
$U$ counts distinct output SHA-256 hashes: identical output bytes count
once for diversity and scanning, even when produced by multiple runs.
A shared hash retains membership in every relevant experimental condition.
Consequently, per-condition unique counts need not sum to the global count.
Time is the recorded duration of a run, including local processing.
Estimated API costs include failed calls and restarted runs. Source records,
version mappings, and the detailed accounting are supplied with the artifact.
